% Ad Familiares 16.17 --- headnote
% ad-familiares-16-17, 29 July 45 BC, Astura

\textit{Cicero from his seaside villa at Astura to Tiro, who is again
unwell and recuperating in Rome. The Perseus dateline is iv K. Sext.
709 (29 July 45 BC), the summer after Tullia's death, when Cicero is
shuttling between Astura and the Tusculan villa and writing at a
furious rate.}

\textit{The letter has two moods stacked on top of each other. The
first section is a teasing literary scolding: Tiro has slipped into a
loose use of \textit{fideliter} (``faithfully'') in his last note ---
the kind of usage that Cicero's own stylistic conscience, normally
embodied by Tiro himself, would have flagged. Cicero rehearses
Theophrastus's doctrine of the ``modest metaphor'' (\textit{verecunda
tralatio}) and promises to take it up in person. The second section
turns practical: Demetrius is arriving, Cicero himself is moving on,
and underneath everything sits the real anxiety --- Tiro's health.
The closing ``count yourself as being with me, if you take care of
yourself'' is the same affectionate calculus that runs through the
whole Tiro cluster.}
